% Section 2.3, from lectures/L02/README.md: what you should be able to do, the questions to test
% yourself with, and the next lecture.

\needspace{12\baselineskip}
\section{Sammanfattning}
\label{c2:sec:summary}

\needspace{6\baselineskip}
\subsection*{Det här ska du kunna nu}
Efter det här kapitlet ska du kunna:
\begin{itemize}
  \item Rita symbolen för varje grundgrind, i ANSI- och i IEC-form, och ange dess sanningstabell och
    uttryck.
  \item Skriva ett booleskt uttryck med bokens notation, och förenkla det med räknelagarna och De
    Morgans lagar.
  \item Ta fram ett uttryck i SP-form ur en sanningstabell, och realisera det som ett grindnät.
  \item Analysera ett grindnät: ta fram dess uttryck och sanningstabell med hjälp av
    mellansignaler.
  \item Bygga och simulera ett grindnät i CircuitVerse, och kontrollera det mot en sanningstabell.
  \item Översätta mellan ett uttryck och ett kontaktnät: serie är AND, parallell är OR, brytande
    kontakt är NOT.
  \item Förklara hur ett relä med självhållning kan minnas att det har startats.
\end{itemize}

\needspace{6\baselineskip}
\subsection*{Frågor att testa dig själv med}
\begin{enumerate}
  \item Vilken grind ger \code{1} bara när alla ingångar är \code{1}? Vilken ger \code{0} bara när
    alla ingångar är \code{0}?
  \item Varför är \mbox{\code{A + A'B}} lika med \mbox{\code{A + B}}? Visa det både med en
    sanningstabell och med ord.
  \item Hur flyttar De Morgan en inversion genom en grind, och varför betyder det att en NAND-grind
    kan bygga vilken funktion som helst?
  \item I bromsassistenten läggs förarens pedal in med OR sist, i stället för att gå genom
    felhanteringslogiken. Vad går sönder om du kopplar den tvärtom?
  \item Varför behöver en XOR-funktion med tryckknappar både den slutande och den brytande
    kontakten i varje knapp?
  \item Vad händer i självhållningskretsen om stoppknappen av misstag kopplas med sin slutande
    kontakt i stället för sin brytande?
\end{enumerate}

CircuitVerse, \url{https://circuitverse.org/simulator}, är simulatorn som grindnäten byggs och
testas i, i det här kapitlet och fram till och med kapitel~\ref{c6:ch}. Innan Labb~1: läs
\secref{c2:app:b} om kontaktnät, och gör förberedelseuppgifterna i bilaga~\ref{app:lab1}.

\needspace{6\baselineskip}
\subsection*{Nästa kapitel}
Kapitel~\ref{c3:ch} hör ihop med Labb~1, kopplingsnät:
\begin{itemize}
  \item \textbf{Labb~1 -- Kopplingsnät}: logiska funktioner byggda med två tryckknappar och en lampa
    på 24~V.
  \item Att analysera ett kontaktnät man inte har ritat själv, och att hitta felet när lampan inte
    gör som den ska.
  \item Förberedelseuppgifterna i labbhandledningen ska vara klara när du börjar labben.
\end{itemize}
